\section{Sample variance and standard deviation}

Theoretical variance \(\Var(X)\) measures spread around the theoretical mean
\(\E[X]\). For observed data, we need an empirical measure of spread around the
sample mean \(\bar{x}\).

\begin{definitionbox}
For a dataset \(x_1,
\ldots,x_n\) with sample mean \(\bar{x}\), the
\textbf{empirical variance} is
\[
v_n=\frac1n\sum_{i=1}^n (x_i-\bar{x})^2.
\]
The corresponding \textbf{empirical standard deviation} is
\[
\sqrt{v_n}.
\]
\end{definitionbox}

This definition treats the empirical distribution as the object being described.
Each observed value has empirical mass \(1/n\), so the variance of the empirical
distribution is exactly
\[
\frac1n\sum_{i=1}^n (x_i-\bar{x})^2.
\]

\subsection*{A small example}

Consider the data
\[
2,\quad 4,\quad 4,\quad 7,\quad 9.
\]
The sample mean is
\[
\bar{x}=5.2.
\]
The deviations from the mean are
\[
-3.2,\quad -1.2,\quad -1.2,\quad 1.8,\quad 3.8.
\]
The squared deviations are
\[
10.24,\quad 1.44,\quad 1.44,\quad 3.24,\quad 14.44.
\]
Therefore
\[
v_n=\frac{10.24+1.44+1.44+3.24+14.44}{5}=\frac{30.8}{5}=6.16.
\]
The empirical standard deviation is
\[
\sqrt{6.16}\approx 2.48.
\]
This number gives a typical scale of spread around the sample mean.

\subsection*{The computational formula}

As in theoretical probability, the empirical variance can also be written as
\[
v_n=\frac1n\sum_{i=1}^n x_i^2-\bar{x}^{\,2}.
\]
This formula is useful for computation, but the definition
\[
v_n=\frac1n\sum_{i=1}^n (x_i-\bar{x})^2
\]
explains the meaning: empirical variance is the average squared distance from
the sample mean.

\subsection*{The denominator \(n-1\)}

Many statistics books and software tools use another quantity:
\[
s^2=\frac1{n-1}\sum_{i=1}^n (x_i-\bar{x})^2.
\]
This is often called the \textbf{sample variance}. Its square root \(s\) is often
called the \textbf{sample standard deviation}.

The difference between \(v_n\) and \(s^2\) is the denominator. The denominator
\(n\) describes the spread of the empirical distribution itself. The denominator
\(n-1\) appears naturally in statistical inference, where the goal is to estimate
the unknown variance of a population from a random sample.

At this descriptive stage, the important point is not to pretend that the two
formulas are the same. They answer closely related but different questions:
\begin{center}
\begin{tabular}{lll}
\toprule
Quantity & Formula & Main role \\
\midrule
Empirical variance & \(\frac1n\sum (x_i-\bar{x})^2\) & describes this dataset \\
Corrected sample variance & \(\frac1{n-1}\sum (x_i-\bar{x})^2\) & used for inference \\
\bottomrule
\end{tabular}
\end{center}

The corrected version will become more important later, when sample statistics
are treated as random variables and their behaviour under repeated sampling is
studied.

\begin{remarkbox}
When using software, always check which convention is being used. Some functions
divide by \(n\), while others divide by \(n-1\).
\end{remarkbox}

\subsection*{Sample variance versus theoretical variance}

The theoretical variance
\[
\Var(X)
\]
belongs to a probability model. It is computed using the distribution of \(X\).
The empirical variance
\[
v_n
\]
is computed from observed values.

If a sample is generated from a model with variance \(\sigma^2\), then empirical
or corrected sample variance may be close to \(\sigma^2\) for large samples under
suitable assumptions. But for a finite sample, equality should not be expected.

\subsection*{What spread means in data}

A small empirical standard deviation means that the observed values are clustered
near their sample mean. A large empirical standard deviation means that the
observed values are more spread out. The scale must be interpreted in the units
of the data.

For example, a standard deviation of \(2\) minutes may be large for a process
that usually lasts \(5\) minutes, but small for a process that usually lasts
\(300\) minutes. Spread is meaningful only relative to context.

\begin{summarybox}
Empirical variance measures spread in the observed data. The denominator \(n\)
gives the variance of the empirical distribution; the denominator \(n-1\) gives a
corrected sample variance used later for inference. Both are data-based
counterparts of the theoretical variance \(\Var(X)\), not the same object as
\(\Var(X)\) itself.
\end{summarybox}
